\section{Discussion}
\label{sec:discussion}

%==========================================================================
% \subsection{Emergence of eCRF Effects through Hierarchical Competition}
\label{sec:disc_emergence}
%==========================================================================
% \textbf{Emergence of eCRF effects through Hierarchical Competition.} \commentcd{ESTE PARRAFO LO BORRARIA} The present study elucidates the complex circuit-level mechanisms underlying robust eCRF phenomena. Historically, the source of surround suppression has been heavily debated, with proposed mechanisms ranging from local, intra-areal horizontal connections to purely top-down inter-areal feedback. Our large-scale spiking network model systematically disentangles these contributions, demonstrating that local lateral connections within V1 are functionally insufficient to reproduce the full magnitude of contextual modulation, typically stalling at roughly 30--35\% suppression. We show that high-magnitude surround suppression---reaching approximately 75\%, which tightly aligns with \textit{in vivo} electrophysiological recordings from primate V1---cannot be accounted for by non-specific linear feedback either.

\textbf{Emergence of eCRF effects through Hierarchical Competition.} Our results reveal that profound contextual modulation is an emergent property of hierarchical competitive dynamics. It is only when top-down feedback signals are non-linearly refined via lateral, winner-take-all competition within higher-order visual areas (such as V2) that the model reproduces physiological suppression depths. These findings fundamentally challenge the traditional view of visual processing as a strictly local feature-extraction pipeline. By embedding a biologically plausible laminar microcircuit within a hierarchical competitive framework, our model bridges the theoretical gap between anatomical constraints---such as the specific targeting of L2/3 and L5 by feedback projections---and the complex, global dynamics of sensory integration.

%==========================================================================
% \subsection{Computational Advantages of Non-linear Feedback}
\label{sec:disc_computational}
%==========================================================================
\textbf{Computational Advantages of Non-linear Feedback.} From a computational perspective, embedding lateral competition within the higher-order feedback loop provides an efficient mechanism for sensory filtering and predictive coding. Simple additive or multiplicative top-down modulations can lead to runaway excitation or generalized gain changes that lack spatial and feature specificity. In contrast, our simulations indicate that inter-areal competition acts as a selective spatial filter. When a visual stimulus extends beyond the classical receptive field, the competitive interactions in V2 actively sculpt the resulting feedback to V1, suppressing neural activity that encodes predictable spatial continuities.

This architecture serves two computational functions. First, it prevents the saturation of primary sensory areas during high-contrast, large-scale visual stimulation, maintaining the network in a stable asynchronous-irregular regime \citep{Brunel2000}. Second, it enhances information transmission by silencing neural populations that encode redundant contextual information, allowing the network to prioritize salient features such as boundaries, discontinuities, or unexpected stimuli \citep{Vinje2000}. This top-down competitive mechanism can therefore be understood as a circuit-level operation for sparse coding and sensory optimization, consistent with predictive coding frameworks \citep{Rao1999, Friston2005}.

%==========================================================================
% TODO: SUBSECCIÓN NUEVA - Implications for Machine Learning
%==========================================================================
% [AQUÍ FALTA: Esta es la subsección MÁS IMPORTANTE para NeurIPS. Debes
% explicitar cómo tu hallazgo puede traducirse en algoritmos o arquitecturas
% para ML. Debes cubrir:
%   - Principio general: inhibición lateral en capas superiores como
%     mecanismo de normalización dependiente de contexto
%   - Analogía con mecanismos existentes en deep learning:
%       * Attention (self-attention, cross-attention)
%       * Squeeze-and-excitation networks
%       * Contrastive learning (SimCLR, MoCo)
%       * Divisive normalization en redes convolucionales
%   - Cómo tu arquitectura podría implementarse en SNNs para tareas de
%     clasificación o segmentación
%   - Ventajas potenciales: robustez adversarial, eficiencia energética,
%     procesamiento relativo vs absoluto
%
%``\subsection{Implications for Machine Learning}
\label{sec:disc_ml}
\textbf{Implications for Machine Learning.} The hierarchical competitive mechanism we identify offers a concrete architectural principle for the design of artificial neural networks. The core insight, that lateral inhibition in higher processing stages enables robust, context-dependent normalization, aligns with several recent trends in deep learning yet remains underexplored as an explicit design choice.

First, our winner-take-all gating of feedback bears functional similarities to attention mechanisms \citep{Desimone1995}, where queries selectively weight the relevance of different spatial locations or feature channels. However, whereas standard self-attention computes pairwise similarities across all positions, incurring quadratic complexity, our lateral inhibition operates locally between modules, offering a potentially more scalable alternative for spatial normalization. This principle aligns with recent work on squeeze-and-excitation networks \citep{Carandini2011}, which recalibrate channel-wise features via a lightweight gating mechanism, and could be extended to incorporate lateral competition between feature modules.

Second, the normalization dynamics we observe, in which V1 responses are scaled relative to the activity of competing V2 modules, implement a form of contrastive computation \citep{Oord2018, Chen2020}. Our biologically constrained architecture suggests that hierarchical lateral inhibition could serve as an energy-efficient alternative to explicit contrastive losses, potentially enabling on-chip learning in neuromorphic hardware \citep{Davies2018}.

Third, the dynamic range compression achieved by competitive feedback offers a potential route to adversarial robustness. By normalizing responses relative to spatial context rather than absolute input magnitudes, the network becomes inherently less sensitive to perturbations that uniformly scale input intensities, a property consistent with known features of divisive normalization in biological vision \citep{Carandini2011}. Finally, our work demonstrates that these computations can be realized in fully spiking architectures, which are compatible with emerging neuromorphic platforms such as Intel's Loihi \citep{Davies2018} and SpiNNaker \citep{Furber2014}. Translating the hierarchical competitive motif to trainable spiking networks, for instance via surrogate gradient methods \citep{Neftci2019}, represents a promising direction for developing low-power artificial perception systems that inherit the robustness and efficiency of biological sensory processing.

%==========================================================================
% \subsection{Implications for Computational Psychiatry}
\label{sec:disc_psychiatry}
%==========================================================================
% \textbf{Implications for Computational Psychiatry.} Beyond basic sensory processing, the mechanistic framework proposed here offers insights into the pathophysiology of neuropsychiatric disorders. The delicate balance between excitation and inhibition (E/I balance) required to sustain hierarchical competition in our model is inextricably linked to the generation of high-frequency oscillatory activity. As our results demonstrate, robust contextual suppression is accompanied by prominent gamma-band oscillations ($\sim 75$ Hz). \commentcd{ESTE RESULTADO ESTARA EN EL APENDICE?} WhThird, the dynamic range compression achieved by competitive feedback offers a potential route to adversarial robustness. By normalizing responses relative to spatial context rather than absolute input magnitudes, the network becomes inherently less sensitive to perturbations that uniformly scale input intensities, a property consistent with known features of divisive normalization in biological vision \citep{Carandini2012}. Finally, our work demonstrates that these computations can be realized in fully spiking architectures, which are compatible with emerging neuromorphic platforms such as Intel's Loihi \citep{Davies2018} and SpiNNaker \citep{Furber2014}. Translating the hierarchical competitive motif to trainable spiking networks, for instance via surrogate gradient methods \citep{Neftci2019}, represents a promising direction for developing low-power artificial perception systems that inherit the robustness and efficiency of biological sensory processing.en we systematically perturbed this local E/I balance to simulate pathological endophenotypes---specifically by attenuating superficial inhibitory synaptic efficacy---the network exhibited a severe degradation in surround suppression alongside a distinct spectral shift and reduction in gamma power. This neurodynamic signature closely mirrors the sensory gating deficits and aberrant oscillatory profiles frequently observed in patients with schizophrenia. Consequently, our spiking network model transitions from a purely theoretical construct to a powerful \textit{in silico} translational platform for computational psychiatry. It provides a cohesive, circuit-level explanation for how localized synaptic micro-deficits can propagate through macroscopic cortical hierarchies, ultimately manifesting as well-documented clinical endophenotypes.

%==========================================================================
% \subsection{Limitations and Future Directions}
\label{sec:disc_limitations}
%==========================================================================
\textbf{Limitations and Future Directions.} While our model successfully captures the core non-linear dynamics of eCRF modulation, several avenues remain for future theoretical and empirical investigation. The current implementation relies on computationally efficient, simplified point-neuron models and utilizes generalized inhibitory population pools. However, cortical inhibition is highly diverse. Future iterations will significantly benefit from disentangling specific interneuron subtypes. Furthermore, the present model operates with static synaptic weights. Incorporating spike-timing-dependent plasticity (STDP) or other biologically plausible learning rules could elucidate how these hierarchical competitive networks develop their precise topographic connectivity during visual maturation. On the practical side, currently the computational implementations of the spiking neural networks relay on CPU. Adaptations to use more efficient compute, as GPUs or neuromorphic, should improve the efficiency of this simulation, opening the door for scaling network size.

%==========================================================================
% TODO: PÁRRAFO DE CIERRE - Conclusión breve
%==========================================================================
% [AQUÍ FALTA: Un párrafo final breve (2-3 oraciones) que sintetice el
% mensaje central y le dé un cierre contundente a la discusión. NeurIPS
% valora un cierre claro.
%
In summary, our work demonstrates that lateral inhibitory competition in higher cortical areas explain robust surround suppression in visual areas. By grounding this principle in a large-scale spiking microcircuit, we provide a mechanistic bridge between physiological data and computational principles, offering a tangible blueprint for the next generation of robust, energy-efficient artificial perception systems.